% !TEX root = ../main.tex
% =====================================================================
% 10. SKILLS ACQUIRED
%
% Sections 5-9 cover the work and the achievements. This short section
% covers the third thing the examiner asked for: the skills.
% Rule I follow: name the skill AND the moment in the project where I
% used it. "I learned PyTorch" says nothing on its own.
% =====================================================================

\section{Skills Acquired}
\label{sec:skills}

\subsection{Scientific skills}

\begin{itemize}\itemsep3pt
  \item \textbf{Reading a paper closely enough to rebuild it.} Reproducing \millet{}'s Conjunctive
        pooling and its AOPCR and NDCG@$n$ metrics (\Cref{sec:bg:millet}) meant going past the
        equations and into the authors' code -- and then noticing where my honest re-implementation
        still did not land on their published number, and being able to name the likely cause: the
        training recipe, not the architecture (\Cref{sec:disc:reading}).

  \item \textbf{Designing a controlled comparison.} Quoting a competitor's published number is not
        enough, so I re-trained FCN, ResNet and InceptionTime myself under \seanet{}'s exact recipe
        (\Cref{sec:setup:baselines}). Any difference in \Cref{sec:results} can then only come from
        the architecture. The same habit produced the ablations in \Cref{app:ablations}, where I
        change exactly one line of a config and nothing else.

  \item \textbf{Reading a metric critically instead of trusting it.} The same architecture scores
        AOPCR 2.57 under my recipe and 13.27 under \millet{}'s (\Cref{sec:disc:reading}). AOPCR is
        unnormalised, and once I understood that, a confusing result turned into a measurable one.
        Transferring Gated and Dual-stream pooling from cancer-pathology MIL taught the same lesson
        from the other side: both underperformed (\Cref{sec:disc:negative}), so an idea borrowed
        from another field is a hypothesis, not a shortcut.
\end{itemize}

\subsection{Technical skills}

\begin{itemize}\itemsep3pt
  \item \textbf{Designing modules around a contract, not an implementation.} Every encoder maps
        $(B, C, \nsteps) \to (B, \dmodel, \nsteps)$ and every pooling head maps
        $(B, \dmodel, \nsteps)$ to the same output. That single rule is the only reason 72
        encoder/pooling combinations could be trained without writing 72 model classes.

  \item \textbf{Working on a shared, time-limited HPC platform.} Reserving GPU nodes on Grid'5000
        with OAR, planning around each job's wall-clock limit, and running the same recipe on two
        sites (Lille and Sophia) so a queue on one did not stall the whole sweep. I also wrote the
        job scripts and a small tmux + Telegram setup that pushed each model's accuracy to my phone
        as it finished, so a failed sweep cost me hours instead of a day.

  \item \textbf{Making ``regenerate everything'' a single command.} One command rebuilds every
        figure and table in this report from the stored results (\Cref{app:commands}), which caught
        a real failure: some generated tables predated a later re-run, so a few numbers were quietly
        out of date. The related habit is asking ``where does this number live after the job ends?''
        -- the per-epoch loss curve went only to MLflow on the compute node and never came home, so
        I changed the training code to save it next to the results.
\end{itemize}

\subsection{Professional skills}

\begin{itemize}\itemsep3pt
  \item \textbf{Reporting progress, including bad news.} I met my supervisor at least once a week
        (\Cref{sec:context}). The hardest update was the O2 discrepancy (\Cref{tab:objectives}):
        re-training \millet{}'s own architecture did not reproduce their published accuracy. My
        first instinct was to keep debugging until it matched. Instead I brought the discrepancy
        itself to the meeting, with the exact numbers and the two causes I could think of. That
        turned a private problem into a shared decision -- we emailed the authors and built a
        full-fidelity reproduction to measure the recipe gap -- and it is where I learned that a
        clearly-stated unresolved result is worth more to a supervisor than a delayed clean one.

  \item \textbf{Planning around a resource I did not control.} With a full sweep costing hours per
        model on shared nodes, deciding what to screen on \webtraffic{} first and what to commit to
        the full sweep (\Cref{sec:setup:data}) was a planning decision as much as a technical one.
        \Cref{sec:res:ucr} shows I got part of that trade-off wrong, which is worth knowing.

  \item \textbf{Writing so that someone else can pick this up.} Explaining MIL and AOPCR to a reader
        who has not read the \millet{} paper forced me to check I understood the mechanisms well
        enough to explain them, not just well enough to implement them; and a living project-state
        note plus a docstring on every module mean the next person does not have to re-derive the
        pipeline from the code.
\end{itemize}

\subsection{The one lesson I would keep}

A careful protocol can overturn a conclusion I had already half-believed. Before the 85-dataset
\ucr{} sweep, the \webtraffic{} numbers told a clean story: my heads win on accuracy, AOPCR and
NDCG@$n$ at once, for a fraction of the parameters. Running the same models against the same
re-trained baseline over 85 more datasets (\Cref{sec:res:ucr}) showed the story does not fully hold.
Neither number is wrong; only the second one is complete. That is the concrete version of ``do not
trust a result you only checked once'', and it is why \Cref{sec:disc:limitations} is a real section
of this report and not a formality at the end.
